\documentclass[conference]{IEEEtran}

\usepackage{amsmath,amssymb}
\usepackage{booktabs}
\usepackage{url}
\usepackage[hidelinks]{hyperref}

\title{Design and Verification of an Instrumented\\
Software-Defined Radio Processing Pipeline}

\author{
\IEEEauthorblockN{Tyree Franklin Jr.}
\IEEEauthorblockA{
Computer Science Student\\
Western Governors University}
}

\begin{document}
\maketitle

\begin{abstract}
This paper examines a software-defined radio processing
platform that combines a C++20 signal-processing core,
a Rust process supervisor, and Python analysis tools.
The architecture uses a bounded producer-consumer queue
and an iterative radix-2 fast Fourier transform to process
simulated in-phase and quadrature samples.
Reanalysis of 1,500 archived measurements gives a mean
processing latency of 0.111014 ms and an interpolated
99th percentile of 0.167423 ms. Every logged block identifies
frequency bin 20, consistent with the intended 5 kHz input.
Source inspection also identifies simulator interface and
serialization defects that prevent these historical results
from validating the reviewed revision.
The analysis defines the measurement boundary, explains
the implications of backpressure and process restart,
and specifies acceptance tests for a reproducible baseline.
\end{abstract}

\begin{IEEEkeywords}
software-defined radio, digital signal processing,
concurrency, performance measurement, software verification
\end{IEEEkeywords}

\section{Introduction}

An SDR processing application must coordinate sample
production, buffering, numerical computation, and output.
A fast transform alone does not establish correct system
behavior. Queue delays, data representation, scheduling,
and failure handling can affect the result independently
of the signal-processing algorithm.

The TeleComms--Software-Defined-Comms repository implements
a hardware-independent platform for studying these
interactions \cite{repo}. C++ provides the processing path,
Rust manages the child process, and Python summarizes
recorded measurements.

This paper asks what the implementation and archived data
establish about processing performance and what additional
verification is needed. Its contributions are an analysis
of the software architecture, a recalculation of the archived
statistics, and a test specification tied to observed
implementation risks.

The work is a student engineering case study. It does not
claim a new transform algorithm, flight qualification,
or demonstrated hard real-time behavior.

\section{Evidence and Method}

The reviewed repository revision is:
\begin{center}
{\footnotesize
\texttt{4c38af8eea230fa8f93ecf38ad5f0bb63f0bbbf5}}
\end{center}

The review covers the C++ source, Rust supervisor, Python
analytics, CMake configuration, FFT test, and committed
metrics CSV. The repository was accessed on October 20, 2026.

Source inspection and archived measurements are treated
as separate evidence. Statements about control flow and
representation follow from the inspected code. Numerical
results are recalculated from the committed CSV.
Proposed tests are identified as future work.

No C++ or Rust build, hardware experiment, or fault-injection
campaign was executed for this analysis. The binary IQ
recording was not numerically evaluated. Accordingly,
the archived measurements are not presented as a reproduction
of the reviewed revision.

This approach follows the distinction between inspection
and testing used in NASA's verification guidance
\cite{verification}. Referencing that guidance does not
establish compliance with NASA requirements.

\section{Architecture}

\subsection{Processing Path}

A producer generates blocks of IQ samples and submits them
to a bounded queue. A consumer removes each block, converts
the samples to the FFT representation, computes the magnitude
spectrum, and selects its largest bin.

\begin{table}[t]
\caption{Intended Live-Mode Configuration}
\label{tab:config}
\centering
\begin{tabular}{@{}lr@{}}
\toprule
Parameter & Value \\
\midrule
Nominal sample rate & 1,024,000 samples/s \\
Tone frequency & 5,000 Hz \\
Transform length & 4,096 samples \\
Queue capacity & 16 blocks \\
Run length & 1,500 blocks \\
Noise amplitude parameter & 0.3 \\
\bottomrule
\end{tabular}
\end{table}

The queue uses a mutex, two condition variables, and a
\texttt{std::deque}. A producer waits when the queue is full.
A consumer waits when it is empty. Both waits recheck their
conditions after waking. Shutdown wakes waiting threads
and allows the consumer to drain remaining items.

Backpressure bounds queue occupancy by slowing the producer.
This behavior is suitable for a cooperative simulator.
It does not prove that a physical radio can deliver samples
without loss when the application falls behind.

The queue copies sample vectors on insertion and removal.
These copies require \(O(N)\) work per operation for a block
of \(N\) samples. Queue payload storage is \(O(QN)\) for
capacity \(Q\). The latency recorder separately retains
one observation per block, so its storage grows with run length.

\subsection{Timing Model}

The nominal block period is
\begin{equation}
T_b=\frac{N}{f_s}
=\frac{4096}{1{,}024{,}000}
=4\ \mathrm{ms}.
\end{equation}

The producer generates and enqueues a block before sleeping
for \(T_b\). Generation, copying, blocking, and scheduler
delay therefore add to the interval between releases.
The configured sample rate describes the signal model,
not a measured wall-clock acquisition rate.

A periodic acquisition model would require absolute release
times, recorded release jitter, and a policy for missed
releases.

\subsection{Process Supervision}

The Rust supervisor launches the C++ executable using Tokio.
A nonzero exit or an injected termination enters a retry
path with a one-second delay. Failure injection is selected
with a configured probability of 30\%, with a requested
termination delay between 50 and 299 ms.

The one-second delay is not a recovery measurement.
Recovery also includes process startup and the time required
to produce valid output.

The constant \texttt{MAX\_RESTARTS} is five, but the counter
is incremented before the limit check. Under repeated child
failures, this permits five launches and four relaunches.
The final status message can report six attempts.
The implementation needs one explicit definition of whether
the limit counts launches or retries.

\section{Signal Processing}

\subsection{Intended Input}

The intended complex signal is
\begin{equation}
x[n]=e^{j2\pi f_0n/f_s}+\eta_I[n]+j\eta_Q[n].
\end{equation}

The noise components are generated by scaling
\texttt{rand()} to a bounded interval. The implementation
does not generate Gaussian noise.

For independent uniform components on \([-a,a]\),
the modeled complex noise power is \(2a^2/3\).
With a unit-magnitude tone and \(a=0.3\), the corresponding
analytical signal-to-noise ratio is
\begin{equation}
10\log_{10}\left(\frac{1}{2a^2/3}\right)
\approx12.22\ \mathrm{dB}.
\end{equation}
This value describes the assumed model. It was not measured
from the archived recording.

\subsection{Transform and Peak Selection}

The implementation computes the unnormalized DFT:
\begin{equation}
X[k]=\sum_{n=0}^{N-1}x[n]e^{-j2\pi kn/N}.
\end{equation}

An iterative radix-2 algorithm performs bit reversal followed
by butterfly stages. The transform requires \(O(N\log_2N)\)
arithmetic work. Magnitude calculation and peak selection
each require \(O(N)\) work.

For the intended configuration,
\begin{equation}
\Delta f=\frac{f_s}{N}=250\ \mathrm{Hz},
\qquad
k_0=\frac{f_0}{\Delta f}=20.
\end{equation}

An ideal unit-amplitude complex tone at bin 20 produces
an unnormalized magnitude of 4,096 at that bin.
This magnitude is not calibrated RF power or power spectral
density.

No window is applied. A bin-centered tone is therefore
a favorable test case. Off-bin tones, negative frequencies,
multiple tones, and varying amplitudes are needed to evaluate
broader spectral behavior.

\section{Performance Analysis}

\subsection{Measurement Boundary}

Live-mode timing starts after queue removal. It includes
sample conversion, spectrum computation, and peak search.
The spectrum function takes its argument by value, so the
timed work also includes that argument copy.

The measurement excludes sample generation, queue waiting,
queue removal copies, CSV output, and console output.
The appropriate term is \emph{processing latency}.
It is neither FFT-only execution time nor end-to-end latency.

For observation \(i\),
\begin{equation}
L_i=t_i^{\mathrm{stop}}-t_i^{\mathrm{start}}.
\end{equation}

\subsection{Archived Dataset}

The CSV contains 1,500 observations with unique, contiguous
block identifiers from 0 through 1,499. All observations
were retained.

The mean is
\begin{equation}
\overline{L}=\frac{1}{M}\sum_{i=1}^{M}L_i.
\end{equation}

For sorted observations, the interpolated percentile uses
position \(h=p(M-1)\), interpolating between the adjacent
values. The C++ summary instead selects the observation
at index \(\lfloor0.99(M-1)\rfloor\).
Table~\ref{tab:results} reports both conventions.

\begin{table}[t]
\caption{Recalculated Archived Results}
\label{tab:results}
\centering
\begin{tabular}{@{}lr@{}}
\toprule
Statistic & Value \\
\midrule
Observation count & 1,500 \\
Mean & 0.111014 ms \\
Median & 0.104865 ms \\
Interpolated p99 & 0.167423 ms \\
C++ convention p99 & 0.167409 ms \\
Minimum & 0.088159 ms \\
Maximum & 0.501117 ms \\
Peak bin in every row & 20 \\
Observations above 4 ms & 0 \\
\bottomrule
\end{tabular}
\end{table}

The interpolated p99 is approximately 4.19\% of the nominal
block period. The maximum is approximately 12.53\%.
These values show processing-time margin in the archived
sample, within the stated measurement boundary.

They do not establish a worst-case bound. The dataset lacks
acquisition timestamps, hardware details, scheduler information,
and an identified producing commit.

The peak bin is 20 in every row, consistent with the intended
5 kHz tone. This checks the recorded peak outputs against
the nominal configuration; it does not verify all complex
transform values.

The blocks represent 6,144,000 samples, equivalent to six
seconds of modeled signal time. Elapsed runtime was not
established.

\section{Source Findings}

\subsection{Simulator Interface}

The simulator header declares global types
\texttt{IQSim} and \texttt{IQSample}.
The pipeline expects \texttt{IQSimulator}, while the entry
point also references \texttt{sdr::IQSimulator} and
\texttt{sdr::IQSample}. Matching declarations are absent
from the inspected source.

The simulator constructor also declares frequency before
sample rate, while callers pass sample rate before frequency.
Correcting only the class name would leave this argument
contract inconsistent.

These findings prevent treating the archived CSV as evidence
that the reviewed revision builds and runs successfully.
They are static findings, not results from an executed compiler.

\subsection{Recording Representation}

IQ sample fields are declared as \texttt{double}.
The recorder writes \texttt{sizeof(float)} bytes from each
field's address. This copies part of a double representation;
it does not convert the value to float32.

The reader similarly writes float-sized byte sequences into
double fields. Python interprets the recording as float32.
The three components therefore lack a consistent numerical
serialization contract.

The repair should use explicit float32 conversions and
float32 read temporaries. The format also needs a defined
byte order, sample rate, version, and malformed-input policy.

\subsection{FFT Input Length}

The FFT assumes a power-of-two length without checking it.
Replay can return a partial final block. An unsupported
length can violate butterfly indexing assumptions and lead
to out-of-bounds access.

Replay must reject unsupported lengths or apply a documented
padding policy before calling the transform.

\subsection{Delivery and Recovery}

The consumer increments the delivered counter, but the
inspected pipeline does not invoke the dropped-block counter.
The dashboard displays a literal zero for dropped blocks.

Contiguous CSV identifiers establish completeness of the
logged range. They do not measure upstream sample loss.
Delivery verification requires separate counts for received,
accepted, processed, rejected, and pending blocks.

Restart also does not imply state recovery. The supervisor
does not restore input offsets or checkpoints, and its
ordinary waiting path has no progress timeout for a stalled
child. File-open errors in record and replay modes can return
without a nonzero process exit, obscuring failure from the
supervisor.

\subsection{Test Coverage and Provenance}

The supplied FFT test checks the positive-frequency peak
of a 1,024-point real cosine at bin 32. It is a useful smoke
test, but does not cover numerical error, recording fidelity,
queue shutdown, or recovery.

The README describes a CI directory that is absent from
the reviewed tree. CI execution is not established by this
snapshot. The CSV also lacks compiler, CPU, operating-system,
seed, and producer-revision metadata.

\section{Verification Plan}

The following are proposed acceptance criteria.
They have not been executed as part of this paper.

\subsection{Build Reproduction}

After interface corrections, a clean checkout shall build
and pass its CMake tests. Each benchmark shall record the
commit, compiler, flags, dependencies, operating system,
CPU, command, seed, and output hashes.

\subsection{Numerical Correctness}

Test lengths 64, 256, 1,024, and 4,096 against an independent
double-precision reference. Inputs shall include an impulse,
DC, positive and negative complex tones, off-bin tones,
and seeded random samples.

For nonzero reference output, the proposed tolerance is
\begin{equation}
\frac{\|X_{\mathrm{test}}-X_{\mathrm{ref}}\|_2}
{\|X_{\mathrm{ref}}\|_2}\leq10^{-4}.
\end{equation}
Zero input shall be tested separately. Unsupported lengths
shall be rejected before transform indexing.

\subsection{Recording Fidelity}

Known finite values shall survive recording and replay
with values equal to the specified float32 conversion.
C++ and Python shall decode identical fixtures.

Tests shall cover empty files, truncated sample pairs,
invalid metadata, and I/O failure. Operational failures
shall produce nonzero process exit codes.

\subsection{Queue Behavior}

Tests shall confirm bounded occupancy, ordered processing,
and exactly-once handling of accepted blocks during normal
completion. Shutdown shall wake a producer blocked on a full
queue and a consumer blocked on an empty queue.
Zero capacity shall be rejected.

At normal completion,
\begin{equation}
B_{\mathrm{accepted}}=B_{\mathrm{processed}}.
\end{equation}
Failure cases shall additionally account for discarded
and pending blocks.

\subsection{Timing and Recovery}

Following correctness verification, collect at least ten
independently launched runs per configuration, each with
at least 15,000 blocks. Declare the warm-up policy and retain
the raw observations.

Measure processing latency, queue residence, end-to-end
latency, release jitter, occupancy, and deadline misses.
Include nominal load, CPU contention, and consumer overload.

A proposed nominal target is p99 processing latency below
4 ms. Passing this target would establish an empirical
processing criterion, not a hard real-time guarantee.

Recovery tests shall deterministically exercise normal exit,
nonzero exit, launch failure, forced termination, and stalled
execution. Measure recovery from fault occurrence to the first
valid output of the replacement process. Verify exact launch
counts at retry exhaustion.

\section{Conclusion}

The project combines numerical processing, concurrent
buffering, process supervision, and offline analysis in
a compact SDR-oriented codebase. Its archived data show
0.111014 ms mean processing latency and 0.167423 ms p99
latency, with consistent detection of the intended tone bin.

The reviewed source cannot yet reproduce those claims
without correcting simulator interfaces and IQ serialization.
Those corrections, followed by the specified verification
tests, would establish a stronger basis for performance
comparison and eventual hardware integration.

\section*{Development Disclosure}

The repository identifies AI assistance in the FFT test
and Python analytics. AI assistance was also used in preparing
this manuscript and recalculating the archived statistics.
The author should verify the final text and accept
responsibility for its claims before distribution.

\begin{thebibliography}{00}

\bibitem{repo}
T. Franklin Jr.,
``TeleComms--Software-Defined-Comms,''
GitHub repository, revision
4c38af8eea230fa8f93ecf38ad5f0bb63f0bbbf5,
accessed Oct. 20, 2026. [Online]. Available:
\url{https://github.com/tyreefranklinjr/TeleComms--Software-Defined-Comms/tree/4c38af8eea230fa8f93ecf38ad5f0bb63f0bbbf5}

\bibitem{verification}
NASA, ``SWE-028: Verification Planning,''
NASA Software Engineering Handbook,
accessed Oct. 20, 2026. [Online]. Available:
\url{https://swehb.nasa.gov/spaces/7150/pages/16450576/SWE-028+-+Verification+Planning}

\end{thebibliography}

\end{document}
